\cotefile{132}
\scanpage{19}
% upper diagram and its branches
\diagram{\begin{tikzcd}[ampersand replacement=\&,row sep=1.2em,column sep=.65em]
\& \mathrm{SO} \& \& \& \\
\& \& \mathrm{PLO}\arrow[ul]\arrow[dl,""{name=pl,coordinate}]\arrow[drr,""{name=pr,coordinate}] \& \& \\
\text{clair} \& \unsure{\mathrm{TUBo}|\mathrm{TUB0}|\mathrm{TUBO}}\arrow[dl,"{\unsure{\simeq|\approx}}"']\arrow[dr,""{name=ts,coordinate}] \& \& \& \mathrm{SOB}\arrow[dl,"{-\partial\;(\simeq\partial)}"]\\
\unsure{\mathrm{Rec}|\mathrm{Rcc}|\mathrm{Kc}} \& \& \mathrm{SOB}\arrow[r,"\partial"] \& \mathrm{CO} \& \\
\& \& \& \rotatebox{90}{$\simeq$} \& \\
\& \& \& \unsure{\mathrm{Ens}\,f|\mathrm{Ensf}|\mathrm{Equiv}} \& \\
\& \& \& (\mathbb{N},\{\unsure{\mathfrak{S}_n|\mathcal{S}_n|\mathfrak{G}_n}\}) \& \text{comme,}\\
\& \& \text{il suffit de les déterminer} \& \&
\arrow[from=pl,to=3-1,bend right=18,thin]
\arrow[from=pr,to=7-5,bend left=10,thin]
\arrow[from=ts,to=8-3,bend right=56,thin]
\end{tikzcd}}

% brace to the right of PLO, upper then lower member
\[\left\{\begin{aligned}
 &\text{complexes \unsure{hypercirculoïdisés|hypercircularisés|hypercirculoïsés|hypercirculaires} (i.e.\ 1-complexes \deleted{\unsure{dont}} circuloïdisés \(K\)}\\
 &\quad\text{avec une partie \(\Pi'\) de \(\mathrm{Circ}(K)\) (ens.\ des circuits)) et}\\
 &\quad\text{application \(\pi'\xrightarrow{g}\mathbb{N}\).}\\
 &\text{extensions des groupes de symétries de 1-complexes \unsure{combinat.|combinatoires}}\\[-2pt]
 &\quad\text{hypercirculoïdisés \deleted{\unsure{hyperboliques}}\inserted{spéciaux} par des groupes de Teichmüller.}
\end{aligned}\right.\]

% below PLO / along the two outgoing arrows
(car \(\pi_0(K)\to\pi_0(S)\)) surjectif \deleted{\unsure{\(g\)}}

% upper right
\(\prod_i T'_{g_i,n_i}\) ; ces ext. sont connues quand on connaît les \(T'_{g_i,n_i}\) séparément, \deleted{\unsure{et}} les hom. canoniques
\[
 \mathbb{Z}^{n_i}\unsure{\lhook\joinrel\to|\to|\xrightarrow{\lambda}}\mathrm{Centre}(L_{g_i,n_i})
\]
et les ext. de Ladegaillerie.
\[
 1\longrightarrow T'_{g_i,n_i}\longrightarrow L_{g_i,n_i}
 \longrightarrow\mathfrak{S}_{\deleted{\unsure{k}}n_i}\longrightarrow1
\]

\margin{cf ci-contre}

% left, beneath Rec
\margin{\(\Leftrightarrow\) [1-complexes combinatoires circuloïdisés] pr. \unsure{quelles|quels} les classes d'isomorphie de 1-complexes circuloïdisés connexes, groupes de symétries des complexes circuloïdisés : doit se ramener à la détermination des groupes de symétries des \unsure{sous-complexes 1-circulaires|sous-complexes 1-circularisés|1-complexes circulaires} connexes}

% middle beneath the left SOB
(familles finies de couples \((g_i,n_i)\)) (genre, nb de trous) \(g\geq0,\ n\geq0\) ; ``groupes de Teichmüller à trous'' ; ou : familles finies de groupes extérieurs à lacets

Tout se détermine \unsure{à|:|: à} l'aide des ``gr. de Teichmüller à trous'' ordinaires \(T_{g,n}\)

% right; fine line from the sloping branch. The framed lines are hatched out
\deleted{\unsure{groupes des auto}\inserted{en termes des} \unsure{groupements} \unsure{disant que}
\(L_{g_i,n_i}\unsure{\times|\twoheadrightarrow|\to}T_{h_i,n_i}\simeq L_{g_i,n_i}/\unsure{\mathrm{Im}}({\unsure{\mathbb{Z}}}^{n_i})\)
de noyau \(\mathrm{Im}(\mathbb{Z}^{n_i}\to L_{g_i,n_i})\)}
ce qui est dit dessus

% bottom; fine line from TUBo--SOB
pour les \deleted{\unsure{des}} \unsure{ordinaires|cas} connexes i.e les 1-compl. circ connexes.

Le calcul de \(g,n\) est donné par Yves Il reste, pour un tel \(K\) \deleted{\unsure{déterminer}} de déterminer
\[
 \boxed{\mathrm{Aut}_{\mathrm{circ}}(K)_{\deleted{\unsure{0|*}}}
 \longrightarrow\mathrm{Aut}(\unsure{\mathrm{Tub}(K)|\mathrm{Tub}\lvert K\rvert})\simeq T_{g,n}}
\]
% archive sheet number: 18

